\subsection{研究方法}
依据 AscendNPU IR 官方 VecAdd 示例和方言文档分析中端层级，
并使用 MLIR 21.1.8 构造上游 VecAdd，执行逐层降低和 CPU 运行验证。

\subsection{AscendNPU IR VecAdd 的方言与层级}
官方快速入门\cite{ascendexample}的 VecAdd 使用 i16 元素，并区分 GM 与 UB
地址空间。输入先搬入片上缓冲，再进行向量加法，最后搬回输出；设备入口用
HACC 属性表达。编译后还需要 CANN runtime 注册和启动设备程序。


架构说明\cite{ascendarchitecture}把 HFusion、HIVM、HACC 分为高层融合、
硬件相关计算/搬运/同步、异构执行表达等职责。结合 HIVM 文档\cite{hivmref}，
可以看到片上内存和计算流水信息在特定层被显式保留。
官方快速入门从 HIVM 级程序开始，体现片上缓冲、向量计算和数据搬运的配合。

\begin{table}[htbp]
\centering\small
\begin{tabular}{p{0.22\linewidth}p{0.31\linewidth}p{0.36\linewidth}}
\toprule
比较点 & 上游 CPU 实验 & Ascend 官方文档路线\\
\midrule
高层计算 & Linalg 通用逐元素计算 & HFusion 等较高层表达\\
硬件相关操作 & 普通 MemRef 访存 & HIVM 中的计算、搬运、同步\\
内存空间 & 本机普通内存 & GM、UB 等硬件空间\\
执行接口 & MLIR 生成 C 包装函数 & HACC 入口及 CANN runtime\\
证据 & 本机各层输出与运行检查 & 官方示例和架构文档分析\\
\bottomrule
\end{tabular}
\caption{两种路线的职责与证据差异}
\end{table}

\subsection{可执行的上游 VecAdd 设计}
独立设计的 \texttt{prep/mlir/vecadd.mlir} 使用三个
\texttt{memref<16xi16>} 参数。Linalg 的三个恒等索引映射表示
$c[i]=a[i]+b[i]$，迭代属性 parallel 表明各元素计算互相独立。
本实验将其降低为标量循环。
\par\vspace{0.7\baselineskip}\noindent\begin{minipage}[t]{\linewidth}
\begin{lstlisting}
linalg.generic {
  indexing_maps = [affine_map<(i)->(i)>,
                  affine_map<(i)->(i)>, affine_map<(i)->(i)>],
  iterator_types = ["parallel"]
} ins(%a, %b : memref<16xi16>, memref<16xi16>)
  outs(%c : memref<16xi16>) {
^bb0(%x: i16, %y: i16, %old: i16):
  %sum = arith.addi %x, %y : i16
  linalg.yield %sum : i16
}
\end{lstlisting}
\end{minipage}\par\vspace{0.5\baselineskip}
输入数组设为 $a[i]=i$、$b[i]=1$（$0\leq i<16$），C 驱动逐项检查输出。

\subsection{逐层 lowering 与检查点}
上游 MLIR 文档介绍结构化 SCF 与非结构化 CF 的区分\cite{scfref}。
本次根据工具可用 pass\cite{mlirpasses}执行如下管线，每一步均保存文件：
\par\vspace{0.7\baselineskip}\noindent\begin{minipage}[t]{\linewidth}
\begin{lstlisting}[language=bash]
mlir-opt vecadd.mlir -o 00_linalg.mlir
mlir-opt 00_linalg.mlir --convert-linalg-to-loops -o 01_scf.mlir
mlir-opt 01_scf.mlir --convert-scf-to-cf -o 02_cf.mlir
mlir-opt 02_cf.mlir --convert-to-llvm \
    --reconcile-unrealized-casts -o 03_llvm.mlir
mlir-translate --mlir-to-llvmir 03_llvm.mlir -o 04_llvm.ll
\end{lstlisting}
\end{minipage}\par\vspace{0.5\baselineskip}
各阶段输出由验证脚本保存。

\begin{table}[htbp]
\centering\small
\begin{tabular}{lrl}
\toprule
检查点 & 行数 & 实际观察到的操作方言\\
\midrule
00\_linalg & 12 & Linalg、Arith、Func\\
01\_scf & 15 & SCF、MemRef、Arith、Func\\
02\_cf & 21 & CF、MemRef、Arith、Func\\
03\_llvm & 67 & LLVM\\
\bottomrule
\end{tabular}
\caption{MLIR 降低阶段与操作方言}
\end{table}

首先 generic 被展开为 \texttt{scf.for}，循环体出现两个 memref.load、一次
arith.addi 和一次 memref.store。接着循环被拆成比较、条件分支、循环体和回边，
迭代变量通过块参数传递；结构化循环这个概念不再显式保留。
再向 LLVM 方言转换时，MemRef 描述符中的已分配地址、对齐地址、偏移、长度和步长
被显式展开，生成与 C 调用者衔接的包装函数。
最后 \texttt{mlir-translate} 才输出真正的 LLVM IR。
因此 \texttt{llvm.*} 操作所在的 MLIR 文件，与 LLVM 自身的 \texttt{.ll} 文件
并不是同一种文本格式。

\subsection{末端验证与分析}
C 驱动构造与一维 MemRef 对应的描述符，通过
\texttt{\_mlir\_ciface\_vecadd} 调用降低后的程序。
本次 LLVM verifier 检查通过，Clang 编译链接后实际输出：
\par\vspace{0.7\baselineskip}\noindent\begin{minipage}[t]{\linewidth}
\begin{lstlisting}
1 2 3 4 5 6 7 8 9 10 11 12 13 14 15 16
\end{lstlisting}
\end{minipage}\par\vspace{0.5\baselineskip}
驱动逐个断言全部 16 个元素正确，退出码 0。
版本、命令、stderr 和结果保存于 \path{prep/mlir/results/}。
行数从 12 增至 67，主要来自内存描述符与调用接口的展开。

多级 IR 在高层保留索引映射和计算结构，在中层明确循环与访存，
最终转换为指针操作和调用接口。各阶段的类型、控制流和内存表示随降低同步转换。
